\addtoc{ВВЕДЕНИЕ}
\section*{ВВЕДЕНИЕ}

Задачи, в~основе которых лежат уравнения эллиптического типа, появляются почти
в~каждой инженерной дисциплине. К~ним сводятся расчёты температурных полей в
конструкциях, оценка распределения потенциала в~проводящих средах,
моделирование процессов фильтрации и множество других прикладных
задач~\cite{tihonov1999}. Аналитическое выражение решения удаётся получить
лишь в~ограниченном числе случаев~--- обычно при~существенной идеализации
геометрии и упрощённой форме правой части. Поэтому в~практической работе
приходится опираться на~численные методы: разностные схемы, конечные элементы,
итерационные алгоритмы решения возникающих систем~\cite{samarskii1989}.

С~увеличением размерности сеточной задачи определяющим параметром становится не
столько достижимая точность, сколько объём требуемых вычислений. Решение
системы линейных алгебраических уравнений, к~которой сводится разностная
аппроксимация, при~густой сетке превращается в~основную статью расходов
вычислительного времени. Один из~путей преодоления этого ограничения~--- метод
декомпозиции области: исходная задача заменяется набором меньших по~размеру
подзадач, решаемых на~отдельных подобластях~\cite{toselli2005}. В~идеальной
ситуации часть подзадач выполняется параллельно, что позволяет приблизить общее
время вычислений к~времени решения одной локальной задачи.

Сама идея относится к~XIX веку: первый вариант альтернирующего метода
предложил Г.~А.~Шварц~\cite{schwarz1870}. В~современной литературе устоявшимися
формулировками являются мультипликативная и аддитивная схемы. Мультипликативная
обновляет приближение последовательно подобласть за~подобластью, используя на
шаге~$i$ свежие значения с~уже обработанной подобласти~$i-1$. Аддитивная схема
выполняет шаги на~всех подобластях независимо, исходя из~общего предыдущего
приближения, а в~зонах перекрытия итог усредняется~\cite{quarteroni1999}.

Каждая из~схем имеет собственные сильные стороны. Аддитивная очевидным образом
ложится на~параллельное исполнение, поскольку локальные задачи решаются
независимо~\cite{cai1999}. Мультипликативная требует меньше итераций для
достижения той же точности, однако последовательный порядок обновления
затрудняет её распараллеливание~\cite{ortega1991}. Какой из~подходов окажется
выгоднее на~конкретной задаче, заранее предсказать сложно: ответ зависит от
размера сетки, числа подобластей, ширины зоны перекрытия и характеристик
используемого оборудования.

Целью настоящей работы является реализация и сравнение мультипликативной и
аддитивной схем метода Шварца применительно к~стационарной задаче, описываемой
уравнением энергии в~одно- и двумерной постановках.

Для~её достижения были поставлены следующие задачи.
\begin{enumerate}
	\item Получить разностные соотношения для~одномерной и двумерной задачи,
	      обосновать выбор шаблона и порядка аппроксимации.
	\item Разработать программный код решателя на~языке MATLAB, разделив его на
	      функции согласно~подзадачам (построение матриц, шаг метода, головной цикл).
	\item Подготовить интерактивные приложения для~запуска экспериментов и
	      визуализации результатов.
	\item Исследовать зависимость качества сходимости от~размера области
	      пересечения, числа подобластей и распределения узлов сетки.
	\item Сопоставить аддитивную и мультипликативную схемы по~числу итераций
	      и реальному времени работы.
	\item Проверить идею ускорения за~счёт замены прямого решения на
	      подобласти итерационным методом Зейделя.
\end{enumerate}

Литература по~методам декомпозиции области достаточно обширна. Базовая теория
и условия сходимости мультипликативной и аддитивной схем подробно изложены в
монографии Тоселли и Видлунда~\cite{toselli2005}, а также в~работе Куартерони и
Валли~\cite{quarteroni1999}. В~обоих источниках обоснована роль перекрытия
подобластей: чем больше радиус пересечения, тем быстрее сходимость, а
конкретный вид оценки константы сходимости зависит от~геометрии и порядка
аппроксимации.

Вопросам параллельной организации вычислений посвящена работа
Ортеги~\cite{ortega1991}. Особый интерес для~практической части представляет
обсуждение накладных расходов, возникающих при~дроблении задачи на~большое
число подобластей. Современная трактовка аддитивной схемы как
предобуславливателя приведена у~Кая и Саркиса~\cite{cai1999}; авторы
показывают, что ограничение зоны перекрытия позволяет сохранить качество
сходимости при~существенном снижении объёма пересылок между~процессами.

Поведение алгоритмов при~малой ширине области пересечения изучено Дрейя и
Видлундом~\cite{dryja1994}. В~частности, показано, что слишком узкое
перекрытие приводит к~зависимости числа итераций от~шага сетки, что
существенно ограничивает применимость метода.

Разностные схемы для~эллиптических уравнений, использованные в~работе для
аппроксимации лапласиана, изложены в~\cite{samarskii1989} и~\cite{bahvalov2003}.
Приёмы работы с~плохо обусловленными системами, оказавшиеся полезными при
рассмотрении метода Зейделя как локального решателя на~подобласти, описаны
в~\cite{golub1999}.

\section{Постановка задачи}

\subsection{Краевая задача и разностная аппроксимация}

В~качестве отправной точки выбрана краевая задача для~стационарного
одномерного уравнения с~известной правой частью:
\begin{equation}\label{eq:1d_problem}
	\frac{d^2 u}{d x^2} = f(x),\quad x\in(a,b),
	\qquad u(a)=u_a,\quad u(b)=u_b.
\end{equation}
Здесь~$f(x)$~--- заданная функция,~$u_a$ и~$u_b$~--- значения искомой функции
на~концах отрезка. Постановка~(\ref{eq:1d_problem}) соответствует
одномерному варианту уравнения теплопроводности в~стационарном случае при
нормировке, переводящей коэффициент теплопроводности в~единицу.

Введём на~отрезке~$[a,b]$ сетку~$a=x_0<x_1<\dots<x_n=b$ с
неравномерным шагом~$h_i=x_i-x_{i-1}$. Вторую производную аппроксимируем
симметричной разностью, обеспечивающей второй порядок точности на
неравномерной сетке~\cite{samarskii1989}:
\begin{equation}\label{eq:central_diff}
	\left.\frac{d^2 u}{d x^2}\right|_{x_i}
	\approx \frac{2}{h_i + h_{i+1}}
	\left(\frac{u_{i+1}-u_i}{h_{i+1}} - \frac{u_i-u_{i-1}}{h_i}\right).
\end{equation}
В~частном случае равномерного разбиения~$h_i\equiv h$ выражение упрощается до
известной формы~$u_{xx}\approx (u_{i+1}-2u_i+u_{i-1})/h^2$. Перенос соседних
узлов в~правую часть даёт пару рекуррентных соотношений, удобных при
последовательном обходе подобласти как слева направо, так и справа налево:
\begin{equation}\label{eq:recur}
	\begin{cases}
		u_{i+1} = 2u_i - u_{i-1} + h^2 f_i, \\[8pt]
		u_{i-1} = 2u_i - u_{i+1} + h^2 f_i.
	\end{cases}
\end{equation}

В~обозначениях величины
\begin{equation}\label{eq:coeffs_1d}
	c^{(L)}_i = \frac{2}{h_i\,(h_i+h_{i+1})},\quad
	c^{(R)}_i = \frac{2}{h_{i+1}\,(h_i+h_{i+1})},\quad
	c^{(\mathrm{M})}_i = -\frac{2}{h_i\,h_{i+1}}
\end{equation}
играют роль элементов левой (\lstinline|glob_lower|), правой
(\lstinline|glob_upper|) и главной (\lstinline|glob_main|) диагоналей.
Подстановка~(\ref{eq:central_diff}) в~(\ref{eq:1d_problem}) с~переносом
граничных значений в~правую часть даёт линейную систему
$A\,u_{\mathrm{int}} = \tilde{f}$ относительно~вектора внутренних значений
$u_{\mathrm{int}}=(u_1,\dots,u_{n-1})^\top$. Матрица~$A$ имеет
трёхдиагональную структуру и в~случае равномерной сетки~--- симметрична:
\begin{equation}\label{eq:tridiag}
	A = \begin{pmatrix}
		-2     & 1      & 0      & \cdots & 0      \\
		1      & -2     & 1      & \cdots & 0      \\
		0      & 1      & -2     & \cdots & 0      \\
		\vdots & \vdots & \vdots & \ddots & \vdots \\
		0      & 0      & 0      & \cdots & -2
	\end{pmatrix},\quad
	\tilde{f} = h^2\!
	\begin{pmatrix}
		f_1 - u_a/h^2 \\
		f_2           \\
		\vdots        \\
		f_{n-1} - u_b/h^2
	\end{pmatrix}.
\end{equation}
Прямое решение этой системы используется как эталонное численное приближение,
с~которым в~дальнейшем сопоставляются итерационные результаты метода Шварца.

\subsection{Двумерная постановка}

В~двумерном случае рассматривается уравнение Пуассона на~единичном квадрате
$\Omega=(0,1)^2$:
\begin{equation}\label{eq:2d_problem}
	\Delta u(x,y) =
	\frac{\partial^{2}u}{\partial x^{2}} +
	\frac{\partial^{2}u}{\partial y^{2}} =
	f(x,y),
	\quad (x,y)\in\Omega,
\end{equation}
с~граничными условиями Дирихле на~$\partial\Omega$. В~качестве модельной правой
части выбрана функция~$f(x,y)=\sin(\pi x)\sin(\pi y)$, краевые условия
однородные, за~исключением правой границы, где~$u(1,y)=\sin(\pi y)$.
Аналитическое решение в~такой постановке выражается через~гиперболический
синус:
\begin{equation}\label{eq:exact_2d}
	u(x,y) = \left(\frac{\sinh(\pi x)}{\sinh(\pi)}
	- \frac{1}{2\pi^{2}}\sin(\pi x)\right)\sin(\pi y).
\end{equation}

Введём прямоугольную сетку~$x_0<\dots<x_n$,~$y_0<\dots<y_m$ с~шагами~$h_x$ и
$h_y$. Стандартный пятиточечный шаблон даёт следующую аппроксимацию:
\begin{equation}\label{eq:poisson_disc}
	h_y^2\,u_{i+1,j} + h_y^2\,u_{i-1,j}
	+ h_x^2\,u_{i,j+1} + h_x^2\,u_{i,j-1}
	- 2\bigl(h_x^2+h_y^2\bigr)\,u_{i,j}
	= h_x^2 h_y^2\,f_{i,j}.
\end{equation}
При~равных шагах форма~(\ref{eq:poisson_disc}) сводится к~привычной записи
$u_{i+1,j}+u_{i-1,j}+u_{i,j+1}+u_{i,j-1}-4u_{i,j}=h^2 f_{i,j}$. На~сетках с
существенно различающимися шагами по~осям использование единого значения~$h$
приводит к~значительной потере точности~--- что было подтверждено при
прогонке предварительных тестов.

Разворачивая внутренние значения~$u_{ij}$ в~один вектор по~столбцам,
получаем разреженную систему с~пятидиагональной матрицей порядка
$(n-1)(m-1)$. Граничные значения переносятся в~правую часть.

\subsection{Идея метода Шварца}

Принцип метода~\cite{schwarz1870} состоит в~следующем. Расчётная область
$\Omega$ покрывается семейством пересекающихся подобластей
$\Omega_1,\dots,\Omega_p$, причём~\mbox{$\bigcup_i \Omega_i = \Omega$}. На~каждой
подобласти решается локальная задача того же типа, что и
исходная~(\ref{eq:1d_problem}) или~(\ref{eq:2d_problem}). Условия на
участках~$\partial\Omega_i\setminus\partial\Omega$, лежащих внутри~$\Omega$,
снимаются с~соседних подобластей, в~которых решение к~данному моменту уже
получено или известно с~предыдущей итерации. Процесс повторяется до~тех пор,
пока приближения на~соседних подобластях не согласуются с~заданной точностью.

Для~одномерной постановки на~отрезке~$[a,b]$ типичный вариант декомпозиции
на~две подобласти изображён на~рисунке~\ref{fig:decomposition}. Граничные
точки~$\xi_1$ и~$\xi_2$ задают ширину зоны перекрытия; число узлов сетки,
попадающих в~неё, обозначается далее~$r$ и называется радиусом пересечения.
В~двумерном случае соответствующая картина выглядит как разбиение квадрата
$\Omega$ на~вертикальные полосы~$\Omega_1$ и~$\Omega_2$, общая часть
которых~--- полоса узлов, лежащих между~$\xi_1$ и
$\xi_2$~(рисунок~\ref{fig:decomposition_2d}).

\begin{figure}[h]
	\centering
	\begin{tikzpicture}[x=1cm, y=1cm]
		% Зона пересечения --- штриховка
		\begin{scope}[on background layer]
			\fill[pattern=north east lines, pattern color=gray!70]
			(5,-0.55) rectangle (7,0.55);
		\end{scope}

		% Основной отрезок
		\draw[thick] (0,0) -- (12,0);
		\fill (0,0) circle (2.2pt) node[below=4pt] {$a$};
		\fill (12,0) circle (2.2pt) node[below=4pt] {$b$};

		% Узлы сетки
		\foreach \x in {1,2,...,11} {%
				\draw (\x,-0.08) -- (\x,0.08);
			}

		% Подобласти
		\draw[red, thick, decorate, decoration={brace,amplitude=8pt,raise=2pt}]
		(0,0.18) -- (7,0.18)
		node[midway, above=12pt, red] {$\Omega_1$};
		\draw[blue, thick, decorate, decoration={brace,amplitude=8pt,mirror,raise=2pt}]
		(5,-0.18) -- (12,-0.18)
		node[midway, below=12pt, blue] {$\Omega_2$};

		% Границы зоны перекрытия
		\draw[dashed] (5,-0.55) -- (5,0.85) node[above=2pt] {$\xi_1$};
		\draw[dashed] (7,-0.55) -- (7,0.85) node[above=2pt] {$\xi_2$};

		% Подписи
		\node[gray!60!black] at (6,-1.6) {зона перекрытия};
	\end{tikzpicture}
	\caption{Декомпозиция отрезка~$[a,b]$ на~две пересекающиеся подобласти
		$\Omega_1$ и~$\Omega_2$. Заштрихованная полоса от~$\xi_1$ до~$\xi_2$
		соответствует общим для~обеих подобластей узлам сетки}
	\label{fig:decomposition}
\end{figure}

\begin{figure}[h]
	\centering
	\begin{tikzpicture}[x=0.9cm, y=0.9cm]
		% Зона пересечения --- штриховка
		\begin{scope}[on background layer]
			\fill[pattern=north east lines, pattern color=gray!70]
			(4,0) rectangle (6,5);
		\end{scope}

		% Основной квадрат
		\draw[thick] (0,0) rectangle (10,5);

		% Узлы внутренней сетки
		\foreach \x in {1,2,...,9} {%
				\foreach \y in {1,2,3,4} {%
						\fill (\x,\y) circle (1.0pt);
					}
			}

		% Подобласти --- цветные рамки
		\draw[red, thick] (0.08,0.08) rectangle (5.92,4.92);
		\draw[blue, thick] (4.08,0.08) rectangle (9.92,4.92);

		% Подписи подобластей
		\node[red] at (2.3, 5.4) {$\Omega_1$};
		\node[blue] at (8, -0.5) {$\Omega_2$};

		% Углы основной области
		\node[below=2pt] at (0,0) {$0$};
		\node[below=2pt] at (10,0) {$1$};
		\node[left=2pt] at (0,5) {$1$};

		% Границы зоны перекрытия
		\draw[dashed] (4,-0.05) -- (4,5.4) node[above=2pt] {$\xi_1$};
		\draw[dashed] (6,-0.05) -- (6,5.4) node[above=2pt] {$\xi_2$};

		% Оси
		\draw[->] (0,-0.05) -- (10.8,-0.05) node[right] {$x$};
		\draw[->] (-0.05,0) -- (-0.05,5.6) node[above] {$y$};

		% Подпись
		\vspace{-1cm}
		\node[gray!60!black] at (5,-1.7) {зона перекрытия};
	\end{tikzpicture}
	\caption{Декомпозиция квадрата~$\Omega = (0,1)^2$ на~две вертикальные
		подобласти~$\Omega_1$ и~$\Omega_2$. Заштрихованная полоса соответствует
		узлам сетки, входящим в~обе подобласти одновременно}
	\label{fig:decomposition_2d}
\end{figure}

Различие между~мультипликативной и аддитивной схемами проявляется в~порядке
обновления подобластей. В~мультипликативном варианте обработка подобласти
$\Omega_i$ использует уже пересчитанное на~текущей итерации решение со~всех
$\Omega_j$,~$j<i$. Эта схема является непосредственным аналогом метода
Гаусса---Зейделя для~задачи на~интерфейсе. В~аддитивном варианте все локальные
задачи на~одной итерации решаются исходя из~единого предыдущего приближения,
после~чего значения в~перекрытых узлах усредняются по~числу подобластей,
которым этот узел принадлежит~--- по~аналогии с~методом
Якоби~\cite{ortega1991}.

\subsection{Алгоритмы}

Приведём оба алгоритма в~виде, отражающем фактически реализованный код.
Пусть подобласть~$\Omega_i$ задаётся набором индексов внутренних узлов
$I_i=\{s_i, s_i+1, \dots, e_i\}$, где~$s_i$ и~$e_i$~--- индексы первого и
последнего внутренних узлов подобласти в~глобальной нумерации. Сосед слева
имеет индекс~$s_i-1$ (узел с~уже известным значением), сосед справа~---
индекс~$e_i+1$. Через~$A_i$ обозначим блок матрицы~$A$, соответствующий
индексам~$I_i$; через~$f_i$~--- сужение правой части~$\tilde{f}$ на~эти
индексы. Параметры~$c^{(L)}_i, c^{(R)}_i$ из~(\ref{eq:coeffs_1d}) задают
вклады левой и правой границ соответственно.

В~мультипликативной схеме обновление подобластей выполняется
последовательно. Пусть~$u^{(k)}\in\mathbb{R}^{n+1}$~--- приближение к~решению
на~$k$-й итерации, причём граничные значения~$u^{(k)}_0=u_a$ и~$u^{(k)}_n=u_b$
фиксированы. Тогда переход к~следующему приближению описывается итерацией:
\begin{equation}\label{eq:mult_method}
	\begin{aligned}
		 & u^{(k+1)} \leftarrow u^{(k)};                                \\
		 & \text{для~} i = 1, 2, \dots, p:                              \\
		 & \quad
		f_i^{\,*} = f_i;\quad
		f_i^{\,*}(s_i) \mathrel{-}= c^{(L)}_{s_i}\,u^{(k+1)}_{s_i - 1};\quad
		f_i^{\,*}(e_i) \mathrel{-}= c^{(R)}_{e_i}\,u^{(k+1)}_{e_i + 1}; \\
		 & \quad
		u^{(k+1)}_{I_i} \leftarrow A_i^{-1}\,f_i^{\,*}.
	\end{aligned}
\end{equation}
Ключевая особенность~--- обращение к~уже обновлённым значениям~$u^{(k+1)}$ на
левой и правой границах подобласти, что и обеспечивает порядок Гаусса---Зейделя
для~интерфейсных значений.

В~аддитивной схеме все локальные задачи решаются исходя из~единого
предыдущего приближения, после~чего значения в~перекрытых узлах усредняются по
числу подобластей, к~которым относится данный узел:
\begin{equation}\label{eq:add_method}
	\begin{aligned}
		 & S \leftarrow 0,\quad W \leftarrow 0;                       \\
		 & \text{(возможно параллельно) для~} i = 1,2,\dots,p:        \\
		 & \quad
		f_i^{\,*} = f_i;\quad
		f_i^{\,*}(s_i) \mathrel{-}= c^{(L)}_{s_i}\,u^{(k)}_{s_i - 1};\quad
		f_i^{\,*}(e_i) \mathrel{-}= c^{(R)}_{e_i}\,u^{(k)}_{e_i + 1}; \\
		 & \quad
		\tilde{u}_i \leftarrow A_i^{-1}\,f_i^{\,*};\quad
		S(I_i) \mathrel{+}= \tilde{u}_i;\quad W(I_i) \mathrel{+}= 1;  \\
		 & \text{затем для~} j\in\bigcup_i I_i:\quad
		u^{(k+1)}_j = S_j / W_j.
	\end{aligned}
\end{equation}
В~неперекрытых узлах~$W_j=1$, и значение совпадает с~локальным решением
$\tilde{u}_i$. В~перекрытых~$W_j\geqslant 2$, и итоговое значение получается
как среднее арифметическое решений на~смежных подобластях.

Критерий остановки в~обеих схемах основан на~максимальном по~модулю
изменении приближения за~итерацию по~всем внутренним узлам:
\begin{equation}\label{eq:stop_crit}
	\varepsilon^{(k)} = \max_{i=1,\dots,n-1}
	\bigl|u^{(k+1)}_i - u^{(k)}_i\bigr|
	\;<\;\varepsilon_{\text{tol}}.
\end{equation}
Алгоритмы~(\ref{eq:mult_method})--(\ref{eq:add_method}) непосредственно
переносятся на~двумерный случай: подобласти соответствуют столбцам сеточной
функции, локальная матрица~$A_i$ строится пятиточечным шаблоном для
выбранного диапазона столбцов, а перенос значений на~границах выполняется по
столбцам~$s_i-1$ и~$e_i+1$. Реализации обеих схем приведены в
листингах~\ref{lst:mult} и~\ref{lst:add}.

\subsection{Модельный пример}

Для~контроля корректности численного решения использовались две тестовые
конфигурации. В~одномерной задаче бралось~$f(x)=\sin(x)$; этой правой части
соответствует точное решение~$u(x)=-\sin(x)$. Рассматривались отрезки
$[0,\pi]$ и~$[0,2\pi]$. В~первом случае краевые условия однородны, во
втором~--- включают одно ненулевое промежуточное значение:
\begin{equation}\label{eq:test_bc}
	\left[\begin{array}{ll}
		\displaystyle[0,2\pi]: &
		\begin{cases}
			u(0)=0,   \\
			u(\pi)=0, \\
			u(2\pi)=0,
		\end{cases}             \\[20pt]
		\displaystyle[0,\pi]:  &
		\begin{cases}
			u(0)=0,                        \\
			u\bigl(\tfrac{\pi}{2}\bigr)=1, \\
			u(\pi)=0.
		\end{cases}
	\end{array}\right.
\end{equation}
Второй вариант представляет дополнительный интерес: внутреннее значение
решения отлично от~нуля, что облегчает выявление ошибок, связанных с
неточным переносом граничных условий между~подобластями.

В~двумерной задаче, как уже отмечалось ранее, использовалось уравнение
Пуассона с~правой частью~$\sin(\pi x)\sin(\pi y)$ и неоднородным граничным
условием~$u(1,y)=\sin(\pi y)$, для~которого аналитическое решение задано
формулой~(\ref{eq:exact_2d}).

\section{Практическая часть}

\subsection{Программный комплекс}

Расчёты выполнялись в~среде MATLAB~\cite{matlab2023}. Исходный код проекта
доступен в~открытом репозитории~\cite{schwartz_github}. Структура кода сознательно
сделана модульной: каждая логически отдельная операция вынесена в
самостоятельную функцию. Это упрощает отладку и поддержку, а также позволяет
переиспользовать общие части. Например, функция построения разреженной
двумерной матрицы по~коэффициентам сетки (\lstinline|build_2d_matrix.m|)
применяется и при~формировании эталонного решателя на~всей области, и при
формировании локальных матриц подобластей в~методе Шварца.

Каталог проекта организован следующим образом:
\begin{itemize}
	\item \lstinline|solvers/|~--- функции построения матриц, шаги
	      мультипликативного и аддитивного методов, прямые решатели для~одно- и
	      двумерной задачи;
	\item \lstinline|scripts/|~--- головные циклы методов, объединяющие
	      инициализацию, основной цикл итераций и сбор истории сходимости;
	\item \lstinline|apps/|~--- два интерактивных приложения, построенных на
	      \mcode{uifigure} и предназначенных для~проведения экспериментов с
	      интерактивной настройкой параметров;
	\item \lstinline|utils/|~--- вспомогательные функции, среди~которых
	      \lstinline|warped_linspace.m| для~построения сетки с~управляемой
	      плотностью узлов.
\end{itemize}

Базовый цикл одномерного метода представлен в~приложении~Г
(листинг~\ref{lst:main}). Указатель \mcode{method} ссылается либо на~функцию
\lstinline|schwartz_mult_step|, либо на~\lstinline|schwartz_add_step|;
переключение между~ними определяется логическим параметром \lstinline|is_mult|.
Это позволяет запускать сравнение двух схем при~полностью идентичных входных
данных.

\subsection{Одномерная задача}

Эксперименты с~одномерной задачей были начаты с~верификации корректности.
Сначала рассматривалась задача на~$[0,\pi]$ с~правой частью~$\sin(x)$ и
однородными краевыми условиями. Эталонное решение, полученное прямой прогонкой
системы~(\ref{eq:tridiag}), сопоставлялось с~приближением, выдаваемым методом
Шварца после~сходимости.

\begin{figure}[h]
	\centering
	\includegraphics[width=\textwidth]{app1d.jpg}
	\caption{Одномерное приложение: слева~--- сравнение точного решения и
		решения по~Шварцу; в~центре~--- история итераций по~подобластям;
		справа~--- графики оценочной и фактической ошибок}
	\label{fig:app1d}
\end{figure}

Интерфейс приложения (рисунок~\ref{fig:app1d}) объединяет три типа графиков.
На~левой панели отображается совпадение точного и приближённого решений после
сходимости. На~центральной~--- ход итераций по~каждой подобласти, причём цвет
графика соответствует её номеру. На~правой~--- сопоставление оценки ошибки
по~критерию~(\ref{eq:stop_crit}) с~фактической, посчитанной относительно
эталона. Расхождение между~ними позволяет судить о~надёжности оценочного
критерия остановки.

Первый результат соответствует ожиданиям: мультипликативная схема сходится
быстрее аддитивной на~той же сетке и при~одинаковых параметрах. Это
согласуется с~известным результатом из~\cite{toselli2005}: аддитивный шаг
интерпретируется как метод Якоби для~значений на~интерфейсе между
подобластями, а мультипликативный~--- как Зейделя для~той же системы.
Отдельный экспериментальный вариант аддитивной схемы с~параллельным запуском
локальных задач средствами \mcode{parfor} оформлен в~файле
\lstinline|schwartz_add_step_cont.m|, однако в~одномерной задаче он не даёт
выигрыша: накладные расходы на~создание параллельных задач превышают
выгоду от~одновременного исполнения. В~основном расчёте 1D используется
последовательный обход подобластей.

\subsection{Размер области пересечения}

Один из~управляющих параметров метода~--- радиус пересечения подобластей~$r$,
выраженный в~числе узлов сетки. Первоначальное предположение состояло в~том,
что разумно брать~$r$ небольшим, экономя на~работе с~дублирующими узлами.
Эксперимент показал, что это представление ошибочно.

\begin{figure}[h]
	\centering
	\includegraphics[width=\textwidth]{optimal_radius.png}
	\caption{Случай~$N=10001$, радиус пересечения~$r=n/4$. Сходимость
		достигается за~5 итераций}
	\label{fig:optimal_radius}
\end{figure}

При~выборе~$r$ равным примерно четверти числа узлов сетки ($r\approx n/4$)
число итераций до~сходимости перестаёт зависеть от~$n$ и в~большинстве
проведённых тестов оказывается на~уровне 4--5
(рисунок~\ref{fig:optimal_radius}). Дальнейшее увеличение~$r$ снижает это
значение ещё сильнее: максимальное ускорение, зафиксированное в
экспериментах, достигается при~$r\approx 3n/4$, когда метод сходится за~две
итерации. Однако такое перекрытие близко к~полному дублированию данных, и
выгоды по~сравнению с~прямым решением задачи на~всей области не наблюдается.

Уменьшение~$r$ даёт качественно противоположную картину. При~\mbox{$r=1$} или~\mbox{$r=2$}
число итераций растёт с~увеличением~$n$, что согласуется с~теоретической
оценкой константы сходимости для~метода Шварца с~малым
перекрытием~\cite{dryja1994}.

\subsection{Распределение узлов сетки}

Если сетка равномерная и общее число узлов невелико, попытка увеличить~$r$
наталкивается на~ограничение: подобласти становятся практически одинаковыми и
смысл декомпозиции теряется. Альтернативный путь~--- сгустить узлы у~концов
отрезка, разредив центральную часть. Тогда фиксированному значению~$r$
соответствует более широкая зона перекрытия, при~том же
общем числе неизвестных. Реализация представлена в~файле
\lstinline|warped_linspace.m| (листинг~\ref{lst:warped}):
\begin{equation}\label{eq:warped}
	\tilde{x}(t) = a + (b-a)\cdot\frac{\operatorname{sgn}(t)\,|t|^{e^{k}}+1}{2},
	\quad t\in[-1,1].
\end{equation}
Параметр~$k$ управляет распределением: при~$k=0$ получается равномерная
сетка ($p=e^0=1$), при~$k>0$ узлы смещаются к~концам, а центральная часть
разреживается. Чем глубже провал плотности в~середине, тем больший
относительный размер по~числу узлов имеет зона перекрытия.

\begin{figure}[h]
	\centering
	\includegraphics[width=\textwidth]{warped_grid.png}
	\caption{Случай~$N=10001$,~$r=2$ при~сгущении сетки к~концам отрезка.
		Сходимость достигается за~9 итераций вместо~нескольких сотен на
		равномерной сетке}
	\label{fig:warped_grid}
\end{figure}

Эффект оказался выраженным. На~сетке со~сгущением узлов у~концов отрезка
при~том же~$r=2$ метод сходится за~9 итераций (рисунок~\ref{fig:warped_grid})
против~более сотни на~равномерной сетке. По~существу получен косвенный путь
расширить эффективное перекрытие, не увеличивая радиус~$r$ в~явном виде:
достаточно сделать сами узлы в~зоне перекрытия более разрежёнными
относительно~концов.

\subsection{Количество подобластей}

\begin{figure}[h]
	\centering
	\includegraphics[width=0.85\textwidth]{five_partitions.png}
	\caption{Разбиение на~5 подобластей. Итераций до~сходимости требуется
		существенно больше, чем при~двух подобластях; выигрыша по~реальному
		времени работы не наблюдается}
	\label{fig:five_partitions}
\end{figure}

Поведение метода при~изменении числа подобластей~$p$ оказалось менее
очевидным, чем ожидалось. В~проведённых экспериментах оптимум по~совокупному
времени работы достигается при~$p=2$, а дальнейшее наращивание до~$p=4$
или~$p=5$ даёт заметное замедление. Причины такого поведения следующие:
\begin{itemize}
	\item с~ростом~$p$ возрастает число итераций до~сходимости, поскольку
	      на~одной итерации информация проходит меньшее расстояние по~сетке;
	\item уменьшение размера локальных задач не компенсирует увеличение их
	      числа~--- особенно с~учётом фиксированных накладных расходов на
	      формирование локальных правых частей и распределение данных в
	      \mcode{parfor};
	\item для~качественного выигрыша от~параллельной аддитивной схемы
	      требуется размер локальной задачи, превышающий некоторый порог; ниже
	      этого порога последовательный код эффективнее.
\end{itemize}
На~рисунке~\ref{fig:five_partitions} приведён ход итераций при~$p=5$:
неустойчивость в~зоне перекрытий выражена сильнее, и количество итераций
до~сходимости заметно превосходит наблюдаемое при~$p=2$.

\subsection{Метод Зейделя как локальный решатель}

Поскольку основная вычислительная нагрузка на~одной итерации приходится на
решение трёхдиагональных систем на~подобластях, представляет интерес замена
прямого решателя итерационным с~ограничением числа шагов. В~работе реализован
метод Зейделя~--- функция \lstinline|solvers/seidel.m| с~параметрами допуска
\mcode{tol} и предельного числа внутренних итераций \lstinline|max_iter|.

Скрипт \lstinline|seidel_schwartz_1d_symmetrical.m| использует метод Зейделя
для~локального решения на~каждой подобласти. Была выдвинута гипотеза
о~том, что на~каждом внешнем шаге метода Шварца достаточно получать
грубое приближение, а его уточнение происходит за~счёт самой итерации
декомпозиции.

\begin{figure}[h]
	\centering
	\includegraphics[width=\textwidth]{seidel_diverge.png}
	\caption{Применение метода Зейделя в~качестве локального решателя.
		При~недостаточной точности внутренних итераций пропадает монотонность
		сходимости, и метод может расходиться}
	\label{fig:seidel_diverge}
\end{figure}

На~практике обнаружился эффект, противоречащий исходной гипотезе. После
нескольких внешних итераций метод выходит на~медленный режим, а при~слишком
слабом задании внутренней точности \mcode{tol} в~зоне перекрытий
наблюдается немонотонная сходимость с~резкими отклонениями приближения
(рисунок~\ref{fig:seidel_diverge}). Выигрыш на~отдельных запусках есть, но
он неустойчив. Согласованное управление двумя точностями~--- внешней и
внутренней~--- оказалось практически неудобным, и в~общем виде идея себя
не оправдала.

\subsection{Двумерная задача}

Перенос алгоритма в~двумерный случай выполнен по~той же схеме: каждая
подобласть представляет собой вертикальную полосу в~плоскости~$(x,y)$, а на
интерфейсе переносятся столбцы граничных значений. Локальная матрица
строится функцией \lstinline|build_2d_matrix.m| в~виде разреженной
пятидиагональной матрицы.

\begin{figure}[h]
	\centering
	\includegraphics[width=\textwidth]{app2d.png}
	\caption{Двумерное приложение. Слева~--- панель параметров,
		в~центре~--- сравнение точного и приближённого решений в~виде
		поверхностей, далее~--- ход итераций и графики ошибки}
	\label{fig:app2d}
\end{figure}

В~качестве модельного примера использовалось уравнение
Пуассона~(\ref{eq:2d_problem}) с~правой частью~$\sin(\pi x)\sin(\pi y)$ и
неоднородным правым граничным условием~$u(1,y)=\sin(\pi y)$, для~которого
аналитическое решение записывается формулой~(\ref{eq:exact_2d}). Поскольку
точное решение известно, оно и использовалось как основной эталон при
оценке погрешности. Дополнительно для~контроля чисто разностной
составляющей ошибки применялся полный решатель \lstinline|solve2d.m|,
который не использует декомпозицию и собирает разреженную систему сразу
для~всей области.

Качественные выводы из~одномерной задачи переносятся и сюда: оптимум по
числу подобластей по-прежнему достигается при~$p=2$, а зависимость
поведения метода от~радиуса перекрытия имеет тот же характер. Существенное
различие появляется при~сравнении схем по~реальному времени работы. На
сетках до~примерно~$300\times 300$ узлов мультипликативная схема сохраняет
преимущество: быстрая сходимость перекрывает накладные расходы. При
большем размере локальные задачи становятся достаточно крупными, чтобы
накладные расходы на~\mcode{parfor} перестали играть решающую роль; в
этом режиме аддитивный вариант начинает превосходить мультипликативный по
времени работы.

\subsection{Сводка результатов}

Если объединить наблюдения в~общий список, картина получается следующей.
\begin{itemize}
	\item Оптимальное число подобластей по~времени работы~---~$p=2$ для
	      обеих размерностей задачи.
	\item Оптимальное перекрытие по~числу итераций~---~$r\approx n/4$, при
	      этом число итераций перестаёт зависеть от~размера сетки.
	\item Сгущение сетки у~концов отрезка даёт ускорение в~десятки раз при
	      том же физическом размере перекрытия.
	\item Мультипликативная схема устойчиво выигрывает в~одномерной задаче и
	      в~двумерной задаче на~сетках умеренного размера.
	\item Аддитивная схема превосходит мультипликативную начиная с~сеток
	      приблизительно~$300\times 300$ узлов в~двумерной задаче за~счёт
	      параллельного запуска локальных решений.
	\item Использование метода Зейделя как локального решателя сокращает
	      время одной итерации, но в~общем случае подрывает устойчивость внешней
	      сходимости.
\end{itemize}

\addtoc{ЗАКЛЮЧЕНИЕ}
\section*{ЗАКЛЮЧЕНИЕ}

В~рамках выполненной работы реализованы и сопоставлены мультипликативная и
аддитивная схемы метода Шварца применительно к~стационарной краевой задаче
в~одно- и двумерной постановке. Полученные экспериментальные результаты
согласуются с~известными теоретическими оценками из~\cite{toselli2005,
	quarteroni1999} и одновременно дают вполне конкретные практические
рекомендации: дробление задачи более чем на~две подобласти оказывается
невыгодным, оптимальное перекрытие соответствует приблизительно~$n/4$
узлам, а сгущение узлов у~концов отрезка позволяет ускорить сходимость без
увеличения общего числа неизвестных. Аддитивная схема становится
конкурентной по~реальному времени лишь при~достаточно больших размерах
локальных задач: на~сетках небольшого размера её выигрыш поглощается
накладными расходами на~параллельный запуск.

Отдельно рассмотрена возможность использования метода Зейделя в~качестве
локального решателя на~подобластях. В~общем виде подход оказался
неустойчивым: ослабление внутренней точности приводит к~потере монотонности
внешней сходимости, метод выходит на~плато либо начинает расходиться.
Возвращение к~этой идее имеет смысл при~использовании более устойчивых
локальных схем, например многосеточных, либо при~адаптивном управлении
внутренней точностью.

В~качестве направлений продолжения работы можно отметить декомпозицию по~оси
большей протяжённости или одновременно по~двум осям, использование
условий третьего рода на~интерфейсе подобластей с~целью повышения
робастности, а также адаптацию метода к~нестационарному уравнению
теплопроводности.
